\documentclass[border=10pt]{standalone}
\usepackage[european, straightvoltages]{circuitikz}
\usepackage{tikz}
\usepackage[T1]{fontenc}
\usepackage[utf8]{inputenc}

% Ustawienia globalne
\ctikzset{
    bipoles/thickness=1.2,
    nodes width=0.1,
    font=\sffamily\small
}

\begin{document}

\begin{circuitikz}[scale=1.1, transform shape]

% ==========================================================
% 1. WTYCZKA SAMOCHODOWA (M15-B) - LEWA STRONA (BARDZO DALEKO)
% ==========================================================
\draw (-12, 0) node[draw, thick, fill=lightgray!20, minimum width=2.5cm, minimum height=14cm, rounded corners] (CONN) {};
\node[rotate=90, font=\bfseries] at (CONN.center) {Złącze M15-B (Hyundai)};

% Definicja współrzędnych pinów wtyczki
\coordinate (P_IGN) at (-10.75, 6);   \node[left] at (-10.8, 6) {Pin 3: IGN (+12V)};
\coordinate (P_BAT) at (-10.75, 2);   \node[left] at (-10.8, 2) {Pin 6: BAT (+12V)};
\coordinate (P_ILL) at (-10.75, -2);  \node[left] at (-10.8, -2) {Pin 9: ILL (+12V)}; % Przesunięte w dół
\coordinate (P_GND) at (-10.75, -4);  \node[left] at (-10.8, -4) {Pin 12: GND};
\coordinate (P_CANH) at (-10.75, -6); \node[left] at (-10.8, -6) {Pin 7: CAN H};
\coordinate (P_CANL) at (-10.75, -7); \node[left] at (-10.8, -7) {Pin 1: CAN L};

% ==========================================================
% 2. ESP32 (CENTRUM) - DUŻY ODSTĘP
% ==========================================================
\draw (4, 0) node[draw, very thick, fill=white, minimum width=6cm, minimum height=12cm] (ESP) {};
\node[above, font=\bfseries\large] at (ESP.north) {ESP32-WROOM-32};

% Piny ESP (Lewa - Wejścia/Zasilanie)
\coordinate (E_VIN) at (1, 5);       \node[right] at (1, 5) {VIN (5V)};
\coordinate (E_GND) at (1, 4);       \node[right] at (1, 4) {GND};
\coordinate (E_ADC) at (1, 2);       \node[right] at (1, 2) {GPIO 34 (ADC)};
\coordinate (E_ILL) at (1, -2);      \node[right] at (1, -2) {GPIO 13 (ILL)}; % Zgodne z P_ILL

% Piny ESP (Prawa - Peryferia)
\coordinate (E_3V3) at (7, 4);       \node[left] at (7, 4) {3V3 Out};
\coordinate (E_SPI) at (7, 2);       \node[left] at (7, 2) {VSPI (18,23,4,2,5)};
\coordinate (E_RX) at (7, -3);       \node[left] at (7, -3) {GPIO 16 (RX)};
\coordinate (E_TX) at (7, -4);       \node[left] at (7, -4) {GPIO 17 (TX)};

% ==========================================================
% 3. ZASILANIE (BUCK CONVERTER) - GÓRA PO LEWEJ
% ==========================================================
\draw (-4, 7) node[draw, fill=white, minimum width=3.5cm, minimum height=2cm] (BUCK) {PRZETWORNICA};
\node at (-4, 7.5) {\small 12V $\to$ 5V (Step-Down)};

% Wejście przetwornicy (z IGN)
\draw (P_IGN) -- (-8, 6) -- (-8, 7) -- (BUCK.west);
\draw (-7, 7) to[pC, l=100$\mu$F/25V, *-] (-7, 5.5) node[ground]{}; % C Input

% Wyjście przetwornicy do ESP
\draw (BUCK.east) -- (0, 7) -- (0, 5) -- (E_VIN);
\draw (-1, 7) to[pC, l=100$\mu$F/10V, *-] (-1, 5.5) node[ground]{}; % C Output

% Masa Przetwornicy
\draw (BUCK.south) -- (-4, 5) node[ground]{};

% Masa ESP
\draw (E_GND) -- (-2, 4) node[ground]{};

% ==========================================================
% 4. WOLTOMIERZ (DZIELNIK NAPIĘCIA) - DUŻO MIEJSCA
% ==========================================================
% Linia od wtyczki do rezystora
\draw (P_BAT) -- (-6, 2) to[R, l=100k (R1), -*] (-3, 2) coordinate(DIV);

% Dzielnik do masy
\draw (DIV) -- (-3, 0.5) to[R, l=20k (R2)] (-3, -0.5) node[ground]{};

% Filtr kondensatorowy
\draw (DIV) -- (-1.5, 2) to[pC, l=1$\mu$F, *-] (-1.5, 0.5) node[ground]{};

% Połączenie do ESP (długa prosta linia)
\draw (DIV) -- (E_ADC);

% ==========================================================
% 5. WEJŚCIE ŚWIATEŁ (ILL)
% ==========================================================
\draw (P_ILL) -- (-6, -2) to[R, l=20k] (-3, -2) coordinate(ILL_NODE);
\draw (ILL_NODE) -- (E_ILL);

% Zener zabezpieczający 3.3V
\draw (ILL_NODE) to[zD, l=3.3V, *-] (-3, -3.5) node[ground]{};

% ==========================================================
% 6. PERYFERIA (LCD i CAN) - PRAWA STRONA (ODSUNIĘTE)
% ==========================================================

% LCD GC9A01 (Wysoko)
\draw (13, 3) node[draw, circle, very thick, minimum size=3.5cm, align=center] (LCD) {\textbf{LCD}\\GC9A01};

% Zasilanie LCD
\draw (E_3V3) -- (8, 4) -- (8, 5) -- (13, 5) -- (LCD.north);
\draw (LCD.east) -- (15, 3) node[ground]{};

% Magistrala SPI
\draw (E_SPI) -- (9, 2) -- (LCD.west);
\node[above, font=\scriptsize] at (10, 2.1) {SCK (18), SDA (23), RES (4), DC (2), CS (5)};

% Moduł CAN (Nisko)
\draw (13, -3.5) node[draw, fill=white, minimum width=3.5cm, minimum height=2.5cm] (CAN) {SN65HVD230};

% Zasilanie CAN
\draw (8, 4) -- (8, -2.5) -- (CAN.north); % 3.3V z tej samej linii co LCD
\draw (CAN.east) -- (15, -3.5) node[ground]{};

% Sygnały CAN (RX/TX)
\draw (E_RX) -- (CAN.west |- E_RX);
\draw (E_TX) -- (CAN.west |- E_TX);

% Wyjście na magistralę (Do wtyczki - Bardzo długa pętla dołem)
\draw (CAN.south) -- (13, -5.5) coordinate(CBUS);
\draw (CBUS) -- (13, -8) -- (-9, -8) -- (-9, -6) -- (P_CANH);
\draw (-9, -8) -- (-9, -7) -- (P_CANL);
\node[above] at (2, -8) {Magistrala CAN (Skrętka - Twisted Pair)};

% Główna masa wtyczki
\draw (P_GND) -- (-10.75, -4) node[ground]{};

\end{circuitikz}

\end{document}